


\subsection{Architectural Overview \& Userspace Bridging}
The module serves as the critical translation layer between the
high-level multithreaded server logic and the physical hardware of the
Xilinx ZedBoard. Because Linux strictly separates userspace applications
from physical hardware memory, this module acts as a secure bridge
utilizing standard POSIX file operations. By hiding complex
bit-shifting, SPI polling loops, and timeout safeguards behind simple
functions like \texttt{TempRead()} and \texttt{GetADC()}, it allows the
higher-level multithreaded architecture to remain clean, fast, and
entirely focused on timing and network transmission. The
\texttt{sensor\_read.c} file is a robust, fault-tolerant module.

\begin{minted}[bgcolor = lightgray, tabsize = 4, breaklines]{c}
    #include <fcntl.h>
    #include <unistd.h>
    #include "measdev.h"
    #include "sensor_read.h"
    MeasObj_struct TF_Obj_Snd, TF_Obj_Rcv;
\end{minted}


\begin{itemize}
    \item \textbf{Communication Struct (\texttt{MeasObj\_struct})}: The
          application cannot read physical memory addresses directly.
          Instead, it packs commands into \texttt{TF\_Obj\_Snd} (Transmit)
          and pushes them into the file descriptor \texttt{fd} pointing to
          the custom kernel driver (\texttt{/dev/meascdd}). The hardware
          response is then retrieved via \texttt{TF\_Obj\_Rcv}.

    \item \textbf{Modularity}: By hiding complex bit-shifting and SPI
          polling loops inside this module, the main sensor thread remains
          clean, lightweight, and hardware-agnostic.
\end{itemize}

\subsubsection{Low-Level SPI Protocol Management}
Interacting with the MAX31865 temperature chip requires precise,
timing-aware SPI bus transactions. The functions \texttt{MaxRead()} and
\texttt{MaxWrite()} are engineered to handle the physical delays of
hardware communication safely using strict polling loops and timeout
safeguards.

\subsubsection*{The MaxRead Function}
\begin{minted}[bgcolor = lightgray, tabsize = 4, breaklines]{c}
unsigned int MaxRead(int fd, unsigned int mreg)
{
    unsigned int xdata;
    int k;

    mreg <<= 8;
    TF_Obj_Snd.rnum = 1;
    TF_Obj_Snd.rvalue = mreg;
    write(fd, &TF_Obj_Snd, MEASOBJ_SIZE);
\end{minted}

\subsubsection*{Command Preparation \& Transfer: }
The target register address is shifted left by 8 bits (\texttt{mreg <<= 8}) and pushed to the driver (\texttt{rnum = 1}).
\begin{minted}[bgcolor = lightgray, tabsize = 4, breaklines]{c}
    k = 0;
    xdata = 0xc0000000;
    do {
        TF_Obj_Snd.rnum = REGNUM_ID;
        TF_Obj_Snd.rvalue = 1;
        write(fd, &TF_Obj_Snd, MEASOBJ_SIZE);
        k++;
        read(fd, &TF_Obj_Rcv, MEASOBJ_SIZE); // Read response from LKM
        xdata = TF_Obj_Rcv.rvalue;
    } while (((xdata & 0x01) == 0) && (k < 1000));
\end{minted}

\subsubsection*{The Polling State Machine \& Timeout: }
The code enters a strict \texttt{do-while} loop, querying the driver
for a status update. It checks whether the hardware is ready by
inspecting the lowest bit (\texttt{xdata \& 0x01}). The loop is capped
at \texttt{k < 1000} iterations to prevent system freezes if the
physical wire is disconnected.

\begin{minted}[bgcolor = lightgray, tabsize = 4, breaklines]{c}
    if (k >= 1000) {
        printf(" *** MaxRead timeout.");
    } else {
        TF_Obj_Snd.rnum = REGNUM_ID;
        TF_Obj_Snd.rvalue = 2;
        write(fd, &TF_Obj_Snd, MEASOBJ_SIZE);
        read(fd, &TF_Obj_Rcv, MEASOBJ_SIZE); // Read response from LKM
        xdata = TF_Obj_Rcv.rvalue;
        xdata &= 0x0ff;
    }
    return xdata;
}
\end{minted}

\subsubsection*{Data Extraction: }
Once ready, it reads the final value and masks it (xdata \&= 0x0ff) to ensure a clean 8-bit payload.

\subsubsection{The MaxWrite Function}
To signal a write operation over SPI, the highest bit of the register
address is set to \texttt{1} (\texttt{mreg |= 0x080}).

\begin{minted}[bgcolor = lightgray, tabsize = 4, breaklines]{c}
void MaxWrite(int fd, unsigned int mreg, unsigned int msend)
{
    unsigned int transm_data, xdata;
    int k;

    mreg |= 0x080;
    transm_data = (mreg << 8) | (msend & 0x0ff);
    TF_Obj_Snd.rnum = 1;
    TF_Obj_Snd.rvalue = transm_data;
    write(fd, &TF_Obj_Snd, MEASOBJ_SIZE);
\end{minted}

\subsubsection*{Payload Construction: }
The 16-bit payload combines the register address and the data byte:
\texttt{(mreg << 8) | (msend \& 0x0ff)}. Identical to \texttt{MaxRead()},
it uses the 1000-iteration polling loop to guarantee the write command
was physically clocked into the sensor.

\begin{minted}[bgcolor = lightgray, tabsize = 4, breaklines]{c}
     k = 0;
    xdata = 0xc0000000;
    do {
        TF_Obj_Snd.rnum = REGNUM_ID;
        TF_Obj_Snd.rvalue = 1;
        write(fd, &TF_Obj_Snd, MEASOBJ_SIZE);
        k++;
        read(fd, &TF_Obj_Rcv, MEASOBJ_SIZE); // Read response from LKM
        xdata = TF_Obj_Rcv.rvalue;
    } while (((xdata & 0x01) == 0) && (k < 1000));
    
    if (k >= 1000) {
        printf(" *** MaxWrite timeout.");
    }
}
\end{minted}

\subsubsection*{MAX31865 Temperature Sensor API}
This section wraps the complex \texttt{MaxRead()} and \texttt{MaxWrite()}
logic into three simple lifecycle functions tailored specifically for the
MAX31865 RTD-to-Digital converter.

\paragraph{Initialization (TempPowerUp): }
Writes the hex value \texttt{0xC1} to the configuration register. This turns on VBIAS, enables Auto-Conversion mode, and enables the 50Hz hardware noise rejection filter. A calculated \texttt{usleep(75000)} allows the analog circuitry 75ms to settle.
\begin{minted}[bgcolor = lightgray, tabsize = 4, breaklines]{c}
int TempPowerUp(int fd)
{
    ssize_t ret;

    TF_Obj_Snd.rnum   = 0;
    TF_Obj_Snd.rvalue = 0x022;
    ret = write(fd, &TF_Obj_Snd, MEASOBJ_SIZE);
    if (ret < 0) { perror("Sensor Initialization Failed"); return -1; }

    MaxWrite(fd, 0x00, 0xC1);   // VBIAS=ON, AutoConv=ON, 2-wire, 50Hz
    usleep(75 * 1000);           // 10ms VBIAS settle + 65ms first conversion
    printf("MAX31865 powered up in auto-conversion mode\n");
    return 1;
}
\end{minted}



\paragraph{Data Retrieval (TempRead): }
Fetches the Most Significant Byte (MSB) and Least Significant Byte (LSB). It evaluates \texttt{if (lsb \& 0x01)} to detect hardware faults (like a broken wire) to prevent broadcasting false data. If valid, the 16-bit integer is reassembled and shifted right by 1 (\texttt{>> 1}) to discard the fault bit, isolating the true 15-bit ADC count.
\begin{minted}[bgcolor = lightgray, tabsize = 4, breaklines]{c}
int TempRead(int fd, int *output_raw)
{
    unsigned int msb = MaxRead(fd, 0x01);
    unsigned int lsb = MaxRead(fd, 0x02);

    if (lsb & 0x01) {
        printf("MAX31865 fault: 0x%02X\n", MaxRead(fd, 0x07));
        return -1;
    }

    *output_raw = (int)((((msb & 0xFF) << 8) | (lsb & 0xFF)) >> 1);
    return 1;
}
\end{minted}


\paragraph{Safe Termination (TempShutdown): }
Writes \texttt{0x01} to disable VBIAS and save power during system shutdown.
\begin{minted}[bgcolor = lightgray, tabsize = 4, breaklines]{c}
void TempShutdown(int fd)
{
    MaxWrite(fd, 0x00, 0x01);
    printf("MAX31865 powered down\n");
}
\end{minted}

\subsubsection{Dual-Channel Analog-to-Digital (ADC) Processing}
The GetADC function communicates with the ZedBoard's ADC peripherals, handling two separate voltage channels simultaneously.
\begin{minted}[bgcolor = lightgray, tabsize = 4, breaklines]{c}
void GetADC(int fd, unsigned int *adc0_out, unsigned int *adc1_out)
{
    int  ccount;
    unsigned int xstatus, adc0, adc1;
    TF_Obj_Snd.rnum = 4;        // start conversion
    TF_Obj_Snd.rvalue = 0;
    write(fd, &TF_Obj_Snd, MEASOBJ_SIZE);
    ccount = 0;
    
    do {
        ccount++;
        read(fd, &TF_Obj_Rcv, MEASOBJ_SIZE);        // Read response from LKM
        xstatus = TF_Obj_Rcv.rvalue;
    } while ((xstatus == 0) && (ccount < 10000000));
\end{minted}

\paragraph{Trigger and Wait: }
Writes a \texttt{0} to \texttt{rnum = 4} to signal the FPGA to sample analog voltages. An extended polling loop (\texttt{ccount < 10000000}) provides the hardware ample time to charge its internal capacitors and complete the conversion.
\begin{minted}[bgcolor = lightgray, tabsize = 4, breaklines]{c}
       TF_Obj_Snd.rnum = REGNUM_ID;
    TF_Obj_Snd.rvalue = 5;      // ADC status register
    write(fd, &TF_Obj_Snd, MEASOBJ_SIZE);
    read(fd, &TF_Obj_Rcv, MEASOBJ_SIZE);        // Read response from LKM
    
    adc1 = TF_Obj_Rcv.rvalue;
    adc0 = adc1 & 0x0fff;
    adc1 >>= 16;
    *adc0_out = adc0;
    *adc1_out = adc1;
}
\end{minted}

\paragraph{Bitwise Separation: }
The hardware returns a highly optimized 32-bit package containing both readings.
\begin{itemize}
    \item \textbf{ADC0 Extraction:} Applies a bitwise AND mask (\texttt{adc1 \& 0x0fff}) to isolate the lowest 12 bits (values 0--4095).
    \item \textbf{ADC1 Extraction:} Applies a bitwise right-shift (\texttt{adc1 >>= 16}) to push the upper 12 bits down to the bottom, safely separating the second reading.
\end{itemize}

\subsubsection{Digital I/O (Switches, Buttons, and LEDs)}
This section maps the physical user interface of the ZedBoard directly into the software.

\paragraph{Switch Acquisition (ReadSwitch): }
Queries the main status register. Applying an \texttt{\& 0xff} (Binary: \texttt{1111 1111}) mask to the 32-bit return value instantly zeroes out irrelevant upper bits, isolating the exact ON/OFF states of the 8 physical sliding switches.

\begin{minted}[bgcolor = lightgray, tabsize = 4, breaklines]{c}
    unsigned int ReadSwitch(int fd)
{
    TF_Obj_Snd.rnum = REGNUM_ID;
    TF_Obj_Snd.rvalue = 0;
    write(fd, &TF_Obj_Snd, MEASOBJ_SIZE);
    read(fd, &TF_Obj_Rcv, MEASOBJ_SIZE);
    
    return (TF_Obj_Rcv.rvalue) & 0xff;
}
\end{minted}

\subsubsection*{Push Button Acquisition (ReadButton): }
Queries the exact same register but shifts the result right by 8 bits (\texttt{>> 8}). Because there are only 5 push buttons, it applies a strict \texttt{\& 0x1f} (Binary: \texttt{0001 1111}) mask to isolate just those 5 states.
\begin{minted}[bgcolor=lightgray, tabsize=4, breaklines]{c}
unsigned int ReadButton(int fd)
{
    TF_Obj_Snd.rnum = REGNUM_ID;
    TF_Obj_Snd.rvalue = 0;
    write(fd, &TF_Obj_Snd, MEASOBJ_SIZE);
    read(fd, &TF_Obj_Rcv, MEASOBJ_SIZE);

    return ((TF_Obj_Rcv.rvalue) >> 8) & 0x1f;
}
\end{minted}

\subsubsection*{LED Control (LedWrite): }
Pushes a standard 8-bit hex payload to \texttt{rnum = 0}. The FPGA hardware interprets this payload and instantly illuminates the corresponding sequence of the 8 onboard yellow LEDs, providing real-time visual feedback.

\begin{minted}[bgcolor=lightgray, tabsize=4, breaklines]{c}
void LedWrite(int fd, uint8_t LED_hex_Value)
{
    ssize_t ret;
    TF_Obj_Snd.rnum   = 0;       // write reg0 (lower 8 bits = uzLED)
    TF_Obj_Snd.rvalue = (unsigned int)LED_hex_Value;
    ret = write(fd, &TF_Obj_Snd, MEASOBJ_SIZE);
    if (ret < 0) perror("LED writing failed");
}
\end{minted}








\clearpage 